\textbf{1.} On factorise :
\[
n^3 - n = n(n^2 - 1) = (n-1)\, n\, (n+1).
\]
Il s'agit du produit de trois entiers consécutifs. Parmi deux entiers
consécutifs, l'un est pair, donc le produit est divisible par $2$.
Parmi trois entiers consécutifs, l'un est divisible par $3$, donc le
produit est divisible par $3$. Comme $2$ et $3$ sont premiers entre eux,
le produit est divisible par $2 \times 3 = 6$.

\medskip

\textbf{2.} On écrit :
\[
n^3 + 5n = (n^3 - n) + 6n.
\]
D'après la première question, $n^3 - n$ est divisible par $6$, et $6n$
l'est évidemment. La somme de deux multiples de $6$ est un multiple de
$6$, ce qui conclut. \hfill $\blacksquare$
